% GENERATED FILE. Edit canonical lessons, then run npm run generate:books.
% canonical-source-hash: fnv1a64:b6c2c85cf9f62464
% canonical-lessons: PA-C258-udhar, PA-C258-karza, PA-C258-sikka, PA-C258-not, PA-C258-kamauna

\chapter{Borrowing and Earning}
\label{ch:pa-a2-borrowing-and-earning}
\hlchaptermodality{258}

\begin{chapteropening}
\textbf{By the end of this chapter:} \emph{I can talk about money I owe and money I earn.}

Complete the last lesson of chapter 258: I can talk about money I owe and money I earn.
\end{chapteropening}

\section[\texorpdfstring{\pa{ਉਧਾਰ} (udhār)}{ਉਧਾਰ (udhār)}]{\pa{ਉਧਾਰ} (udhār) — credit; money lent}

\label{lesson:PA-C258-udhar}



\begin{quote}
Before the new one: say the Punjabi for to withdraw, then the Punjabi for to go up.
\end{quote}

\subsection*{You'll want to know: \pa{ਉਧਾਰ}}
\textbf{\pa{ਉਧਾਰ}} — \emph{udhār} — \textquotedblleft{}credit; money lent\textquotedblright{}.

\begin{quote}
\pa{ਮੈਂ} \pa{ਉਧਾਰ} \pa{ਨਹੀਂ} \pa{ਲੈਂਦਾ} \pa{ਹਾਂ।}
\end{quote}

\begin{quote}
\emph{maiṁ udhār nahīṁ laindā hāṁ.}
\end{quote}

\begin{quote}
I do not take things on credit.
\end{quote}

A masculine noun.

\begin{cousinweb}
\textbf{\pa{ਉਧਾਰ}} (\emph{udhār}) is inherited from Sanskrit \emph{uddhāra}, \textquotedblleft{}raising; a loan\textquotedblright{}; CDIAL lists Punjabi \emph{udhār}, \textquotedblleft{}loan\textquotedblright{}, under it (CDIAL 2018).
\end{cousinweb}

\begin{grammarlens}[title={What you've built}]
One more word for reading notices and messages.
\end{grammarlens}

\subsection*{Your turn}
\begin{itemize}
\raggedright
  \item \emph{Say it:} \emph{udhār}
  \item \emph{Say it:} \emph{udhār}, once more
  \item \emph{Say it:} \emph{maiṁ udhār nahīṁ laindā hāṁ.} — I do not take things on credit.
  \item \emph{Recall:} say the Punjabi for to withdraw, then the Punjabi for to go up, then say \emph{udhār} again
\end{itemize}

\begin{tcolorbox}[breakable,colback=teal!4,colframe=teal!35!black,title={Before you move on}]
What does \pa{ਉਧਾਰ} mean? (\textquotedblleft{}Credit; money lent\textquotedblright{}.) Say it once more.
\end{tcolorbox}
\section[\texorpdfstring{\pa{ਕਰਜ਼ਾ} (karzā)}{ਕਰਜ਼ਾ (karzā)}]{\pa{ਕਰਜ਼ਾ} (karzā) — a debt}

\label{lesson:PA-C258-karza}



\begin{quote}
Before the new one: say the Punjabi for to go up, then the Punjabi for credit.
\end{quote}

\subsection*{You'll want to know: \pa{ਕਰਜ਼ਾ}}
\textbf{\pa{ਕਰਜ਼ਾ}} — \emph{karzā} — \textquotedblleft{}a debt\textquotedblright{}.

\begin{quote}
\pa{ਅਮਨ} \pa{ਉੱਤੇ} \pa{ਬੈਂਕ} \pa{ਦਾ} \pa{ਕਰਜ਼ਾ} \pa{ਹੈ।}
\end{quote}

\begin{quote}
\emph{Aman utte baiṅk dā karzā hai.}
\end{quote}

\begin{quote}
Aman has a debt to the bank.
\end{quote}

A masculine noun.

\begin{cousinweb}
\textbf{\pa{ਕਰਜ਼ਾ}} (\emph{karzā}) is a loan, usually traced through Persian to Arabic \emph{qard}, \textquotedblleft{}a loan\textquotedblright{}.
\end{cousinweb}

\begin{grammarlens}[title={What you've built}]
One more word for reading notices and messages.
\end{grammarlens}

\subsection*{Your turn}
\begin{itemize}
\raggedright
  \item \emph{Say it:} \emph{karzā}
  \item \emph{Say it:} \emph{karzā}, once more
  \item \emph{Say it:} \emph{Aman utte baiṅk dā karzā hai.} — Aman has a debt to the bank.
  \item \emph{Recall:} say the Punjabi for to go up, then the Punjabi for credit, then say \emph{karzā} again
  \item \emph{Write it:} copy \textbf{\pa{ਕਰਜ਼ਾ}} once, checking the whole word
\end{itemize}

\begin{tcolorbox}[breakable,colback=teal!4,colframe=teal!35!black,title={Before you move on}]
What does \pa{ਕਰਜ਼ਾ} mean? (\textquotedblleft{}A debt\textquotedblright{}.) Say it once more.
\end{tcolorbox}
\section[\texorpdfstring{\pa{ਸਿੱਕਾ} (sikkā)}{ਸਿੱਕਾ (sikkā)}]{\pa{ਸਿੱਕਾ} (sikkā) — a coin}

\label{lesson:PA-C258-sikka}



\begin{quote}
Before the new one: say the Punjabi for credit, then the Punjabi for a debt.
\end{quote}

\subsection*{You'll want to know: \pa{ਸਿੱਕਾ}}
\textbf{\pa{ਸਿੱਕਾ}} — \emph{sikkā} — \textquotedblleft{}a coin\textquotedblright{}.

\begin{quote}
\pa{ਮੇਰੀ} \pa{ਜੇਬ} \pa{ਵਿੱਚ} \pa{ਇੱਕ} \pa{ਪੁਰਾਣਾ} \pa{ਸਿੱਕਾ} \pa{ਹੈ।}
\end{quote}

\begin{quote}
\emph{merī jeb vich ikk purāṇā sikkā hai.}
\end{quote}

\begin{quote}
There is an old coin in my pocket.
\end{quote}

A masculine noun.

\begin{cousinweb}
\textbf{\pa{ਸਿੱਕਾ}} (\emph{sikkā}) is a loan, usually traced through Persian to Arabic \emph{sikka}, \textquotedblleft{}a die for stamping coins\textquotedblright{}.
\end{cousinweb}

\begin{grammarlens}[title={What you've built}]
One more word for reading notices and messages.
\end{grammarlens}

\subsection*{Your turn}
\begin{itemize}
\raggedright
  \item \emph{Say it:} \emph{sikkā}
  \item \emph{Say it:} \emph{sikkā}, once more
  \item \emph{Say it:} \emph{merī jeb vich ikk purāṇā sikkā hai.} — There is an old coin in my pocket.
  \item \emph{Recall:} say the Punjabi for credit, then the Punjabi for a debt, then say \emph{sikkā} again
  \item \emph{Write it:} copy \textbf{\pa{ਸਿੱਕਾ}} once, checking the whole word
\end{itemize}

\begin{tcolorbox}[breakable,colback=teal!4,colframe=teal!35!black,title={Before you move on}]
What does \pa{ਸਿੱਕਾ} mean? (\textquotedblleft{}A coin\textquotedblright{}.) Say it once more.
\end{tcolorbox}
\section[\texorpdfstring{\pa{ਨੋਟ} (noṭ)}{ਨੋਟ (noṭ)}]{\pa{ਨੋਟ} (noṭ) — a banknote}

\label{lesson:PA-C258-not}



\begin{quote}
Before the new one: say the Punjabi for a debt, then the Punjabi for a coin.
\end{quote}

\subsection*{You'll want to know: \pa{ਨੋਟ}}
\textbf{\pa{ਨੋਟ}} — \emph{noṭ} — \textquotedblleft{}a banknote\textquotedblright{}.

\begin{quote}
\pa{ਮੇਰੇ} \pa{ਕੋਲ} \pa{ਸੌ} \pa{ਦਾ} \pa{ਨੋਟ} \pa{ਹੈ।}
\end{quote}

\begin{quote}
\emph{mere kol sau dā noṭ hai.}
\end{quote}

\begin{quote}
I have a note for one hundred.
\end{quote}

A masculine noun.

\begin{cousinweb}
\textbf{\pa{ਨੋਟ}} (\emph{noṭ}) is borrowed from English \emph{note}.
\end{cousinweb}

\begin{grammarlens}[title={What you've built}]
One more word for reading notices and messages.
\end{grammarlens}

\subsection*{Your turn}
\begin{itemize}
\raggedright
  \item \emph{Say it:} \emph{noṭ}
  \item \emph{Say it:} \emph{noṭ}, once more
  \item \emph{Say it:} \emph{mere kol sau dā noṭ hai.} — I have a note for one hundred.
  \item \emph{Recall:} say the Punjabi for a debt, then the Punjabi for a coin, then say \emph{noṭ} again
  \item \emph{Write it:} copy \textbf{\pa{ਨੋਟ}} once, checking the whole word
\end{itemize}

\begin{tcolorbox}[breakable,colback=teal!4,colframe=teal!35!black,title={Before you move on}]
What does \pa{ਨੋਟ} mean? (\textquotedblleft{}A banknote\textquotedblright{}.) Say it once more.
\end{tcolorbox}
\section[\texorpdfstring{\pa{ਕਮਾਉਣਾ} (kamāuṇā)}{ਕਮਾਉਣਾ (kamāuṇā)}]{\pa{ਕਮਾਉਣਾ} (kamāuṇā) — to earn}

\label{lesson:PA-C258-kamauna}



\begin{quote}
Before the new one: say the Punjabi for a coin, then the Punjabi for a banknote.
\end{quote}

\subsection*{You'll want to know: \pa{ਕਮਾਉਣਾ}}
\textbf{\pa{ਕਮਾਉਣਾ}} — \emph{kamāuṇā} — \textquotedblleft{}to earn\textquotedblright{}.

\begin{quote}
\pa{ਮੇਰੀ} \pa{ਭੈਣ} \pa{ਦਫ਼ਤਰ} \pa{ਵਿੱਚ} \pa{ਚੰਗੇ} \pa{ਪੈਸੇ} \pa{ਕਮਾਉਂਦੀ} \pa{ਹੈ।}
\end{quote}

\begin{quote}
\emph{merī bhaiṇ daftar vich chaṅge paise kamāundī hai.}
\end{quote}

\begin{quote}
My sister earns good money at the office.
\end{quote}

Past (I …): \pa{ਮੈਂ} \pa{ਬਹੁਤ} \pa{ਪੈਸੇ} \pa{ਕਮਾਏ}, \emph{maiṁ bahut paise kamāe}.

\begin{cousinweb}
\textbf{\pa{ਕਮਾਉਣਾ}} (\emph{kamāuṇā}) is inherited from a reconstructed \emph{*karmāpayati}, \textquotedblleft{}works, earns\textquotedblright{}; CDIAL lists Punjabi \emph{kamāuṇā}, \textquotedblleft{}to work, earn\textquotedblright{}, under it (CDIAL 2897).
\end{cousinweb}

\begin{grammarlens}[title={What you've built}]
One more word for reading notices and messages.
\end{grammarlens}

\subsection*{Your turn}
\begin{itemize}
\raggedright
  \item \emph{Say it:} \emph{kamāuṇā}
  \item \emph{Say it:} \emph{kamāuṇā}, once more
  \item \emph{Say it:} \emph{merī bhaiṇ daftar vich chaṅge paise kamāundī hai.} — My sister earns good money at the office.
  \item \emph{Recall:} say the Punjabi for a coin, then the Punjabi for a banknote, then say \emph{kamāuṇā} again
  \item \emph{Write it:} copy \textbf{\pa{ਕਮਾਉਣਾ}} once, checking the whole word
\end{itemize}

\begin{tcolorbox}[breakable,colback=teal!4,colframe=teal!35!black,title={Before you move on}]
What does \pa{ਕਮਾਉਣਾ} mean? (\textquotedblleft{}To earn\textquotedblright{}.) Say it once more.
\end{tcolorbox}
